% =============================================================================
% Structured abstract. Compiles via the canonical wrapper; see BUILD.md.
% =============================================================================

\begin{abstract}

\noindent\textbf{Objective.} To test whether stratified within-cohort analysis, bidirectional external validation, and case-level paired bootstrap inference jointly surface failure-mode magnitudes in cross-continental clinical prediction that aggregate metrics conceal.

\vspace{0.5em}
\noindent\textbf{Materials and Methods.} Eight machine learning models (XGBoost and logistic regression, on preoperative and preoperative$+$intraoperative feature sets) were trained on each of INSPIRE (Korea; $n =$~\pnTableOneNInspire) and MOVER (USA; $n =$~\pnTableOneNMover), then evaluated bidirectionally between cohorts. Case-level paired bootstrap (2{,}000 iterations) was the primary inferential framework. Direction asymmetry was stress-tested via matched-subsampling across four case-mix dimensions (ASA, Elixhauser comorbidity, emergency proportion, temporal period). Feature-importance transferability used SHAP rank correlation; calibration used slope, intercept, O:E, and Brier score before and after Platt scaling.

\vspace{0.5em}
\noindent\textbf{Results.} Across the eight cross-population runs, aggregate AUCs concealed substantially lower within-stratum AUCs (Simpson's paradox gap range 5.0--16.5 pp; worst case: aggregate AUC 0.756 vs within-stratum AUCs 0.58--0.60). Cross-continental transferability was markedly direction-asymmetric ($+$8.53 pp, 95\% CI 6.91--10.24, bootstrap $p = 0.001$); the asymmetry survived matching on ASA and comorbidity, attenuated to 70--86\% of baseline after matching emergency proportion, and was untestable for temporal period. Pre-Platt calibration slopes ranged 0.41--1.29 across all eight cross-population runs; 5-fold CV Platt scaling restored slopes to 0.95--1.02. Intraoperative features conferred a mean external AUC advantage of \pnMeanIntraopAdvantageFull.

\vspace{0.5em}
\noindent\textbf{Discussion.} These magnitudes are clinically material and not visible to conventional aggregate reporting. The methodological commitments that surface them are well-established individually; their joint application here characterizes failure-mode magnitudes that single commitments would underestimate.

\vspace{0.5em}
\noindent\textbf{Conclusion.} We present this case study as a cautionary reference for cross-population deployment of clinical prediction models, with reproducibility infrastructure released for verification and extension.

\end{abstract}

% =============================================================================
% END OF ABSTRACT
% =============================================================================
